El presente informe aborda el problema de encontrar diferencias entre dos cadenas de texto, a partir de la subsecuencia común más larga (LCS, por sus siglas en inglés), con el objetivo de comparar el rendimiento de dos paradigmas algorítmicos: una solución por fuerza bruta y una basada en programación dinámica. Se desarrollaron ambas implementaciones considerando una entrada estándar con múltiples casos de prueba. Para evaluar el rendimiento de los algoritmos se generó un conjunto de datos, con 10 casos de prueba diseñados para representar distintos tamaños de entrada, lo que permitió analizar el tiempo de ejecución y uso estimado de memoria en función del tamaño de las cadenas comparadas. Los resultados experimentales muestran una clara superioridad de la programación dinámica frente a la fuerza bruta a medida que el tamaño crece, validando así su idoneidad para este tipo de problemas.